\section{The universal module and sharp metabelian bounds}
\label{sec:universal}

\subsection{A multigraded commutator module}

For \(d\ge2\), put
\[
 R_d=\F_4[t_1,\ldots,t_d]/(t_1^2,\ldots,t_d^2),
 \qquad\m=(t_1,\ldots,t_d).
\]
Let \(W_d\) be the \(R_d\)-module generated by \(v_{ij}\), \(i<j\),
subject to
\begin{equation}\label{eq:module-triangle}
 t_iv_{jk}+t_jv_{ik}+t_kv_{ij}=0\qquad(i<j<k).
\end{equation}
We set \(v_{ji}=v_{ij}\) and \(v_{ii}=0\).
This is the classical free metabelian commutator-module presentation,
base-changed to \(R_d\); compare~\cite[Theorem~6.2]{PapadimaSuciu2004}
and~\cite[Sections~2--3]{KovacsStohr2013}.

Every derived module in Theorem~\ref{thm:correspondence} is a quotient of
\(W_d\): send \(t_i\) to \(P_i\) and \(v_{ij}\) to \(a_{ij}\).
The intrinsic-generation condition makes this homomorphism surjective.
Conversely, every quotient of \(W_d\), with these images, gives the
normal form \eqref{eq:normal-form}.

Grade \(R_d\) by coefficient degree, and put the generators \(v_{ij}\)
in coefficient degree zero. We also use the finer grading
\(\deg t_i=e_i\), \(\deg v_{ij}=e_i+e_j\).

\begin{proposition}\label{prop:Hilbert}
The coefficient-degree Hilbert polynomial of \(W_d\) is
\begin{equation}\label{eq:Hilbert}
 H_d(T)=\sum_{k=2}^d\binom dk
       \bigl((k-1)T^{k-2}+kT^{k-1}\bigr)+\binom d2T^2.
\end{equation}
Consequently
\[
 \dim_{\F_4}W_d=(d-1)2^d+1-d+\binom d2.
\]
The universal solution on \(\F_4^d\oplus W_d\) has class
\(\max(4,d+1)\).
\end{proposition}
\begin{proof}
A raw multihomogeneous monomial is \(t_Sv_{ij}\), where \(S\) is
square-free. Its multidegree has entries zero, one, or two, and at most
two entries are doubled. Fix its support of size \(k\ge2\).

If no index is doubled, the raw coordinates are the edges of the complete
graph on the support. Relations~\eqref{eq:module-triangle} are its triangle
sums. The quotient edge space has dimension \(k-1\), in coefficient
degree \(k-2\).

If exactly one index \(j\) is doubled, only edges incident with \(j\)
occur. Triangle relations through \(j\) identify every pair of these
edges, leaving dimension one. There are \(k\) choices of \(j\), in
coefficient degree \(k-1\).

If exactly two indices \(i,j\) are doubled, there is one raw coordinate.
A third index in the support kills it by a triangle relation, since the
other two terms contain a squared variable. The coordinate survives
exactly when the support is \(\{i,j\}\), giving the final term of
\eqref{eq:Hilbert}. These cases exhaust the multigrading.

The maximum coefficient degree is two for \(d=2,3\), and \(d-1\) for
\(d\ge4\). Equation~\eqref{eq:lower-central-module} gives the class.
\end{proof}

In particular, binary degree two permits unbounded polar nilpotency class.
For the exact minimum dimension we must retain a terminal commutator in
as small a quotient of \(W_d\) as possible.

\subsection{Terminal tensors and contraction modules}

\begin{lemma}\label{lem:terminal-repeat}
Let \(\mathfrak l\) be a metabelian \(\F_4\)-Lie algebra satisfying
\eqref{eq:Lie-action-square}, of class \(c\ge3\).
Some nonzero terminal commutator has a repeated input between its first
bracket pair and an operator position. If \(c\ge5\), there is a
subalgebra of the same class generated by exactly \(c-1\) elements
modulo its derived ideal.
\end{lemma}
\begin{proof}
The last lower central ideal is spanned by
\[
 T(x_1,x_2;x_3,\ldots,x_c)
     =P(x_3)\cdots P(x_c)\beta(x_1,x_2).
\]
Indeed \(W=\Span\beta+\m W\), and multiplication by \(\m^{c-2}\)
kills the second summand. The operator positions are symmetric, and a
repeated input within them gives zero.

Suppose every value \(T(x,y;x,z_4,\ldots,z_c)\) were zero.
Polarizing the repeated input gives
\[
 T(x,y;z,\ldots)=T(z,y;x,\ldots).
\]
Together with symmetry of the first pair and the operator positions,
this makes \(T\) fully symmetric. Jacobi on the first three inputs
then gives \(3T=T=0\), a contradiction.

Choose a nonzero repeated-input value. Its inputs span a subspace
\(U_0\) of the abelianization of dimension at most \(c-1\).
Their chosen lifts generate a subalgebra \(H\) of class \(c\), with
abelianization exactly \(U_0\): all words of length at least two lie
in its derived ideal. Its module is a quotient of \(W_{d_0}\), where
\(d_0=\dim U_0\).
By Proposition~\ref{prop:Hilbert}, class \(c\ge5\) requires
\(d_0\ge c-1\). Thus \(d_0=c-1\).
\end{proof}

For \(d\ge4\), the top coefficient degree of \(W_d\) is \(d-1\).
It has basis \(T_1,\ldots,T_d\), where
\begin{equation}\label{eq:terminal-basis}
 T_j=\Bigl(\prod_{h\ne i}t_h\Bigr)v_{ij}\quad(i\ne j).
\end{equation}
Triangle relations identify the choices of \(i\).

\begin{lemma}\label{lem:top-contraction}
The top component is the \(\GL(d,4)\)-module \(\det(U)\otimes U\).
For every nonzero functional \(\lambda\) on that component, extended
by zero to the other degrees,
\[
 \dim_{\F_4}R_d\lambda=2^d-2.
\]
If a quotient \(M\) of \(W_d\) retains a nonzero top component, then
\(\dim_{\F_4}M\ge2^d-2\), whether or not the quotient is graded.
\end{lemma}
\begin{proof}
Under a diagonal basis change \(A\), \(T_j\) scales by
\(\det(A)A_{jj}\). A transvection \(e_j\mapsto e_j+\alpha e_k\)
sends \(T_j\) to \(T_j+\alpha T_k\), and fixes the other top basis
vectors. To check the formula choose \(i\notin\{j,k\}\) in
\eqref{eq:terminal-basis}; any extra product has a squared variable.
For the remaining vectors an index outside the affected entries can
be chosen since \(d\ge4\).
Together with permutations these transformations prove the module assertion.
They act transitively on nonzero dual vectors, including scalar multiples.

Normalize \(\lambda(T_j)=1\), \(\lambda(T_k)=0\) for \(k\ne j\).
We claim
\begin{equation}\label{eq:annihilator}
 \Ann_{R_d}(\lambda)=\Bigl(\prod_{i\ne j}t_i\Bigr).
\end{equation}
The annihilator is multigraded: contractions by distinct monomials have
distinct resulting dual multidegrees. For a square-free monomial \(t_S\)
not containing \(j\), the remaining multidegree is of the one-doubled-index
type. It survives unless \(S\) contains every other index.
If \(j\in S\), the remaining multidegree has no doubled index.
It survives whenever an index is omitted, as witnessed by an edge
joining that index to \(j\); it vanishes when all indices are used.
These are exactly the two monomials in the ideal on the right of
\eqref{eq:annihilator}. This proves the claim and the dimension formula.

For an arbitrary quotient \(M\), pull back a functional nonzero on the
top image to \(\phi\in W_d^*\). Its highest homogeneous component
\(\phi_{d-1}\) is nonzero. Independent homogeneous contractions of
\(\phi_{d-1}\) lift to independent contractions of \(\phi\), by
comparison of their leading dual degrees. Hence
\[
 \dim M\ge\dim R_d\phi\ge\dim R_d\phi_{d-1}=2^d-2.
\]
This is the usual filtered contraction argument; it does not assume a
graded kernel for \(W_d\to M\).
\end{proof}

\subsection{Exact minima}

\begin{theorem}\label{thm:unrestricted-minimum}
Among metabelian quadratic solutions of polar class \(c\ge2\), the
minimum binary dimension is
\[
 \mu(2)=6,\quad\mu(3)=8,\quad\mu(4)=12,\qquad
 \mu(c)=2^c+2c-6\quad(c\ge5).
\]
\end{theorem}
\begin{proof}
Write \(d=\dim_{\F_4}(V/W)\), \(w=\dim_{\F_4}W\).
Nonabelian objects have \(d\ge2\) by Lemma~\ref{lem:two-generators}.
A nonzero product \(P_1\cdots P_kv\), for commuting square-zero
operators, yields \(2^k\) independent vectors
\[
 v_S=\Bigl(\prod_{i\in S}P_i\Bigr)v\qquad(S\subseteq\{1,\ldots,k\}).
\]
In a relation, choose a term with nonzero coefficient of smallest support
and multiply by its complementary operators. Every other surviving
term repeats an operator and vanishes, while the chosen term does not.
Taking \(k=c-2\) gives \(w\ge2^{c-2}\).
For \(c=2,3,4\) this gives the respective lower bounds \(6,8,12\).

They are attained by cyclic quotients of \(W_2=R_2v_{12}\): use
\(R_2/\m\), \(R_2/(t_2)\), and \(R_2\), respectively, as the
derived module in \eqref{eq:normal-form}.
The first two representatives are
\begin{align*}
 &(a,b,z)\longmapsto(\omega a,\omega b,\omega z+ab),\\
 &(a,b,y,z)\longmapsto
    (\omega a,\omega b,\omega y+az+a^2b,\omega z+ab).
\end{align*}
The class-four representative appears in
\eqref{eq:seven-coordinates} below.

Let \(c\ge5\). Proposition~\ref{prop:Hilbert} gives \(d\ge c-1\).
Lemma~\ref{lem:terminal-repeat} supplies a same-class subalgebra with
\(d_0=c-1\). Its derived module retains the top of \(W_{c-1}\), so
Lemma~\ref{lem:top-contraction} gives
\(w\ge2^{c-1}-2\). Therefore
\[
 \dim_{\F_2}V=2(d+w)\ge2^c+2c-6.
\]

To attain the bound, set \(d=c-1\ge4\), choose a top functional
\(\lambda\) as in Lemma~\ref{lem:top-contraction}, and put
\begin{equation}\label{eq:terminal-quotient}
 M_d=W_d/\{w:\lambda(rw)=0\text{ for every }r\in R_d\}.
\end{equation}
The denominator is an \(R_d\)-submodule. The quotient has dimension
\(2^d-2\) and retains the terminal value, so its solution has class
\(d+1=c\) and binary dimension \(2(d+2^d-2)\).
\end{proof}

The change of formula at class five comes from the top degree of the
universal module. For two and three quotient generators,
\[
 H_2(T)=1+2T+T^2,\qquad H_3(T)=3+8T+6T^2.
\]
In these cases the degree-two double-endpoint terms contribute at the
top, allowing class four even with two quotient generators. For
\(d\ge4\), they lie below the top degree \(d-1\); the terminal component
\(\det(U)\otimes U\) then controls class \(d+1\) uniformly. This is why
the contraction argument yields one formula from class five onward,
while the cyclic quotients of \(W_2\) give the three smaller values.

The extremal quotient \(M_d\) in \eqref{eq:terminal-quotient} has
coefficient Hilbert polynomial
\[
 (d-1)+\sum_{k=1}^{d-1}\binom d{d-1-k}T^k.
\]
For example, the class-five extremizer has derived Hilbert vector
\((3,6,4,1)\) and binary lower central dimensions
\((36,28,22,10,2,0)\).
The contraction construction supplies an explicit finite module
description at every class, without an orbit search.
